% Separate continuation after B.2.13; exact infinite-to-finite comparison.
\subsection{The exact central-cover comparison for the quotient $U$}

Let $K\subseteq F\cap B'$ be any closed normal subgroup
of $B$, and set
\[
 U=B/K,\qquad U^{\mathrm c}=B/\overline{[B,K]},\qquad
 M=K/\overline{[B,K]},\qquad J=F/K.                            \tag{B.S1}
\]
No finiteness of $U$ or $M$ is assumed.  These definitions
apply to the six-coordinate kernel used in B.2.13.  To check
$K\subseteq B'$ there, take the abelian values of the first
and the $s$ coordinates in the target $C_3$ quotient of
$B^{\mathrm{ab}}$.  For source abelian coordinates $(v,w)$
they are $w$ and $v$ modulo three.  Thus both vanish for
$K$, while the source map $F\to B^{\mathrm{ab}}=C_6$
has value $4(v+w)$, which consequently vanishes.

For any of $B,U,U^{\mathrm c}$, a positive framed
automorphism here means a continuous automorphism fixing
$\tau$, together with specified conjugators of $s,r$ in
the actual image of $F$.  Let $\mathscr P(B;K)$ retain
those automorphisms preserving $K$, and let
$\mathscr P(U^{\mathrm c};M)$ retain those preserving $M$.
Both have literal reduction maps into $\mathscr P(U)$.

\begin{result}{B.2.14}{Proved}{An exact central-cover image equality for infinite quotients}
For the actual marked quotients (B.S1),
\[
 \operatorname{im}\bigl(\mathscr P(B;K)\longrightarrow
                              \mathscr P(U)\bigr)
 =
 \operatorname{im}\bigl(\mathscr P(U^{\mathrm c};M)
                        \longrightarrow\mathscr P(U)\bigr).
                                                                  \tag{B.S2}
\]
The equality retains both prescribed $U$-conjugator values.
It holds for infinite $U$ and infinite $M$ as well.
\end{result}

\begin{proof}
A full automorphism preserving $K$ preserves its closed
commutator with $B$.  Reduction gives the inclusion from
left to right, including the actual conjugators.

For the reverse inclusion, fix a positive framed
automorphism $\widetilde\alpha$ of $U^{\mathrm c}$
preserving $M$, and write $(\alpha,\beta,A)$ for its
reduction to $U$.  Let $N$ range over open characteristic
subgroups of $U$ contained in $J\cap U'$.  These are
cofinal because $U$ is finitely generated and
$U/U'=C_6$.  Put
\[
 R_N=U/N=B/K_N,\qquad
 K_N=\operatorname{preimage}_B(N).
\]
Then $K\subseteq K_N\subseteq F\cap B'$, so
$R_N^{\mathrm c}=B/\overline{[B,K_N]}$ is the finite
canonical marked cover in (B.Q1).

Let $\widetilde N$ be the inverse image of $N$ in
$U^{\mathrm c}$.  The canonical quotient maps satisfy
the exact identity
\[
 \ker(U^{\mathrm c}\longrightarrow R_N^{\mathrm c})
       =\overline{[U^{\mathrm c},\widetilde N]}
       =\overline{[B,K_N]}/\overline{[B,K]}.                    \tag{B.S3}
\]
This follows directly by taking the images of the indicated
commutators under $B\to U^{\mathrm c}$; the containment
$K\subseteq K_N$ gives the last quotient.  Since
$\widetilde\alpha$ preserves $M$ and its induced
$\alpha$ preserves characteristic $N$, it preserves
$\widetilde N$ and the subgroup in (B.S3).  It therefore
induces an actual positive framed automorphism of the
finite group $R_N^{\mathrm c}$, preserving its canonical
central kernel and fixing $\tau$.  Both frames are the
literal reductions of the given ones.

The generalized finite-certificate equality in the proof
of B.2.9 supplies a full positive $\tau$-fixed automorphism
of $B$ inducing the indicated action on $R_N$, with the
two prescribed frame values there.  This holds for every
$N$ and for every common characteristic refinement of
finitely many such $N$.

The full set of positive $\tau$-fixed framed automorphisms
of $B$ is compact.  Indeed it is a closed subset of the
compact space $F^2$ of its two conjugators: impose the
peripheral product relation and, in every finite quotient
of $B$, generation by the two proposed torsion images.
Each condition is closed.  The resulting source map is
onto every finite quotient, hence onto $B$; the Hopfian
property of the finitely generated profinite group $B$
makes it an automorphism.  The conditions of inducing
$(\alpha,\beta,A)$ modulo each $N$ are also closed and
have the finite intersection property by the preceding
paragraph.  Compactness supplies one full framed
automorphism satisfying all of them.  Cofinality identifies
its induced action and both frames on $U$ with
$(\alpha,\beta,A)$.  In particular it preserves $K$.
This proves the reverse inclusion in (B.S2).
\end{proof}

For the particular quotient of B.2.13, (B.S3) is the
explicit comparison with all the finite canonical covers;
it supplies the cofinal passage absent from the quotient
positivity argument alone.  Establishing that its positive
automorphisms (B.R15) occur in the right side of (B.S2)
is still a separate task.  The result gives an exact image
equality and does not assert that this image is all of
$\mathscr P(U)$.
